\chapter{Use Case Diagram}
\label{ch:usecase}

\section{Purpose}

The use case diagram shows what the actor can ask the system to do. It models
interaction at the level of user intent, without reference to the components
that satisfy it.

\umlfigure{use-case-diagram}{Use case diagram}{fig:usecase}

\section{Actors}

\begin{longtable}{L{0.20\textwidth}L{0.72\textwidth}}
	\caption{Actors}\label{tab:actors}\\
	\toprule
	\textbf{Actor} & \textbf{Description} \\
	\midrule
	\endfirsthead
	\toprule
	\textbf{Actor} & \textbf{Description} \\
	\midrule
	\endhead
	Engineer & The primary and only actor. Works in GENESYS, invokes both
	translations, and reviews their results. The actor is deliberately singular
	because the same person performs both directions. \\
	\bottomrule
\end{longtable}

\section{Use Cases}

\begin{longtable}{L{0.26\textwidth}L{0.66\textwidth}}
	\caption{Use cases}\label{tab:usecases}\\
	\toprule
	\textbf{Use case} & \textbf{Description} \\
	\midrule
	\endfirsthead
	\toprule
	\textbf{Use case} & \textbf{Description} \\
	\midrule
	\endhead

	Import Model from Code & The engineer selects source code and the tool
	produces the corresponding model in GENESYS. Realised by the Code Extractor,
	Abstraction and Model Mapper components of Figure~\ref{fig:block}. \\
	\addlinespace

	Export Model to Code & The engineer selects a model and the tool produces
	source code implementing it. Realised by the Model Mapper, Abstraction and Code
	Generator components of Figure~\ref{fig:block}. \\
	\bottomrule
\end{longtable}

\section{System Boundary}

Both use cases sit inside the \emph{Genesys Translation Tool} boundary. GENESYS
itself and the code repository are outside it, consistent with
Figure~\ref{fig:block}, where both are drawn as external systems.